\section{Historical and computational provenance}
\label{app:provenance}\label{sec:provenance}
The publication introducing a numerical method and the publication
establishing a particular operation count are not necessarily the
same. Table~\ref{tab:provenance} separates the algorithmic
antecedent, published numerical-analysis cost results, and the
model-specific arithmetic counted in this article. In particular,
the historical SVD and conjugate-gradient papers do not prove
our composite smoothness-indexed grid-search theorems.

\begin{sidewaystable*}[p]
\centering
\caption{Algorithm provenance and sources of computational complexity.
Original methods are distinguished from later analyses of their costs;
our order-$d$ PLS specialization is identified explicitly.}
\label{tab:provenance}
\setlength{\tabcolsep}{8pt}
\renewcommand{\arraystretch}{1.30}
\normalsize
\begin{tabularx}{0.94\textheight}{@{}p{0.18\textheight}p{0.31\textheight}X@{}}
\toprule
\textbf{Algorithm family} & \textbf{Original method / historical antecedent}
& \textbf{Complexity source and PLS interpretation}\\
\midrule
Dense linear solves &
Classical elimination; sparsity-aware elimination:
\citet{markowitz1957}. &
Dense $O(n^3)$ arithmetic and $O(n^2)$ storage:
\citet{golubvanloan2013}. General sparse costs depend
on fill and ordering \citep{davis2006}.\\
\addlinespace
Reduced full SVD &
Bidiagonal singular-value algorithm:
\citet{golubkahan1965};
practical SVD: \citet{golubreinsch1970}. &
Dense $r\times n$ SVD: $O(nr^2)$ work and
$O(nr)$ retained basis \citep{golubvanloan2013}.
The per-candidate PLS score counts are derived here.\\
\addlinespace
Dense symmetric eigenanalysis &
QR-family eigensolver: \citet{francis1961};
penalty-orthogonal spline basis: \citet{demmler1975}. &
Standard dense $O(n^3)$ work and
$O(n^2)$ full eigenvector storage
\citep{golubvanloan2013}.\\
\addlinespace
Banded Cholesky &
Cholesky's method as documented by
\citet{benoit1924}; earlier structured
WH computation: \citet{henderson1925}. &
For half-bandwidth $d$, $O(nd^2)$
factorization and $O(nd)$ storage
\citep{golubvanloan2013,anderson1999};
WH-specific work: \citet{weinert2007}.\\
\addlinespace
Selected inverse subset &
Early actuarial trace antecedent:
\citet{spoerl1943}; general sparse
recurrences: \citet{takahashi1973,erisman1975}. &
Selected-inverse recurrences and their computational
dependencies: \citet{erisman1975}; practical
supernodal implementation: \citet{lin2011}.
The stated banded $O(nd^2)$ is counted here.\\
\addlinespace
General sparse direct factorization &
Sparsity-aware elimination:
\citet{markowitz1957}. &
Orderings, fill and algorithm-specific storage:
\citet{davis2006}; generic sparsity does
not imply fixed bandwidth.\\
\addlinespace
Conjugate gradients &
Original CG: \citet{hestenesstiefel1952}. &
Preconditioned convergence: \citet{saad2003}.
The $O(qnd)$ cost of $q$ local-difference
iterations (excluding preconditioner costs)
is derived here.\\
\addlinespace
WH smoothing and GCV &
Early graduation: \citet{bohlmann1899};
\citet{whittaker1923}; WH cross-validation:
\citet{brooks1988}. &
Earlier efficient WH/GCV algorithms:
\citet{weinert2007}; classical spectral
GCV in splines and ridge regression:
\citet{craven1979,golub1979}.\\
\bottomrule
\end{tabularx}

\vspace{8pt}
\begin{minipage}{0.94\textheight}
\small\emph{Interpretation.} Original method and complexity
source are separate claims. Bounds shown for a generic matrix
do not establish an optimal bound for the special difference
operator; our complete-grid workload bounds follow in
Section~\ref{sec:complexity}.
\end{minipage}
\end{sidewaystable*}




\paragraph{Sources that extend rather than replace the exact methods.}
The exact finite-sample $d=2$ Hodrick--Prescott
weights studied by \citet{corneamadeira2017}
and the continuous derivative-penalty spline
construction of \citet{reinsch1967}
are separate antecedents with different
problem formulations or computational routes.
Likewise, modern selected-inverse software
\citep{lin2011} implements a general sparse
factorization-and-recursion strategy rather
than originating the historical
\citet{erisman1975} recurrences. The
random-probe trace estimator of
\citet{hutchinson1990} and the
stochastic Lanczos matrix-function method of
\citet{ubaru2017} produce approximate
diagnostics; their error probabilities and
inner iteration costs should not be silently
identified with the exact traces,
determinants and leverage vectors
costed in this article.

\paragraph{Historical priority within trend smoothing.}
The review by \citet{weinert2007} traces penalized-difference
graduation to \citet{bohlmann1899}, predating
\citet{whittaker1923}; documents the early WH factorization
of \citet{henderson1925}; and identifies \citet{spoerl1943}
as an earlier user of an inverse-trace computation later
associated with \citet{takahashi1973}. That review also identifies \citet{brooks1988}
as an early published WH cross-validation application.
Importantly, \citet[Sec.~1, p.~961]{weinert2007} explicitly
reports $O(n)$ arithmetic for a fixed-order hat-matrix trace
and develops especially efficient WH factorization/GCV
algorithms principally for difference order two. Thus
linear-time fixed-order WH diagnostics are pre-existing
results, not a novelty claim here. We cite these findings as reported
by Weinert, without claiming independent verification
of all nineteenth-century and actuarial archival sources.
